\section{Tokenizers Under Test}\label{sec:methods}

All of these tokenizers ingest the same text, and produce an \nVocabBudget-type vocabulary for each language.
They all tokenize on whitespace, but they differ in what additional information they draw on.
Two of them, the baselines, do not use any additional information, and three of them, the POS-weighted variants, receive POS tags for every word of the training set.
The Lemma2Word tokenizer, in addition, receives a lemma for every word of the training set.

\subsection{Byte-Pair Encoding with a Free Space Marker}\label{sec:freebpe}
The byte-pair encoding (BPE) algorithm~\cite{sennrich2016bpe} constructs a vocabulary by repeatedly merging the most frequent symbol pair.
Our implementation of BPE inserts a marker (\textvisiblespace) between words.
This marker can be merged with the word that precedes it (\texttt{the\textvisiblespace}) or with the word that follows, or become internal to a token that spans two words (\texttt{of\textvisiblespace the\textvisiblespace}).
In SentencePiece~\cite{kudo2018sentencepiece}, \textvisiblespace{} is always prefixed to each word at the start of training, so that tokens can never span word boundaries.

In this evaluation, we distinguish between the former, the free-marker variant (\emph{BPE, free marker}), which is the one that we will use to demonstrate POS-weighted tokenization, and the latter, the word-initial marker variant (\emph{BPE, word-initial marker}).
Vocabularies are made up of the alphabet of the text supplemented with the learned merge operations.

\subsection{Part-of-Speech-Weighted Merges}\label{sec:posbpe}
Universal Dependencies POS tags are associated with every word of the training set, and a transition probability table $P(b \mid a)$ is computed from this tagged text.
The POS-weighted variant replaces the pair count with a sum weighted over its occurrences.
The weight of joining the end of a word with the marker following it, which we are interested in, is $1 + p$, where $p$ is the transition probability between the POS tags of those words, clamped between $0.01$ and $0.95$.

If the marker is bare on either side, the weight is further multiplied by $0.98$, and all other cases have a weight of 1.
This scheme encourages tying a trailing space to its word when the next word is predictable from its tag, slightly discourages interaction with the marker, and has no impact on internal pairs.
A first ablation is to drop the $0.98$ multiplier, which we term \emph{transitions only}.
A second is to drop the transition term, which we call \emph{penalty only}.
In the earlier iteration of this implementation, the scores were recomputed every time they were needed, so they could only be computed for small vocabularies.
The updated implementation maintains a list of all occurrences for each pair, and updates their scores after each merge, which allows training of an \nVocabBudget-type vocabulary on 3,000 paragraphs in less than a minute.

There are two defects in the earlier implementation.
First, transitions that had never been observed were assigned a probability of $0.6$, which is more than most of the observed values; this has been changed to the floor value of $0.01$.
Second, in the older version, the word to which each symbol belonged was recovered by searching for bare markers in the current symbol sequence, so that when a marker was merged it was no longer counted, and later computed weights of transitions applied to a different pair of words.
In the updated version, the word for each symbol is recorded when the symbol sequence is built.
If both of these fixes are disabled, the rewritten code computes the same pair scores before the first merge, and the same unweighted merge sequence, as the original, as can be confirmed in the tests in \texttt{experiments/test\_posbpe.py}.

\subsection{Content-Word Weighting}\label{sec:content}
The goal of the POS signal is to preserve content words (such as \emph{scientist}) intact rather than splitting them into components.
The merge weighting method in Section~\ref{sec:posbpe} does not achieve this goal because it never changes the score of a pair inside a word.
The simplicity of the approach led us to propose a variant that weights all pairs within words having POS tags \textsc{noun}, \textsc{propn}, \textsc{verb}, and \textsc{adj} by two and all other pairs by one.
The weight of two was selected before running the experiments and never tuned.
This variant moves vocabulary capacity toward whole content words and away from function words and affixes.

\subsection{Lemma2Word}\label{sec:l2w}
To develop Lemma2Word, we first obtain the Stanza lemma of each word occurring in the training text.
The stem is the lemma if it is a substring of the word (\emph{walk} in \emph{walking}); otherwise we take the longest common substring of length at least three, such as \emph{habl} for the words \emph{hablo} and \emph{hablar}.
Anything before the stem is a prefix and anything after is a suffix; at most two boundaries can occur in a word.
We derive three items from the training text.
The first item is a lexicon that maps each word to its boundaries, using majority vote to resolve inconsistent lemmatization.
The other two are the inventory of suffixes and of prefixes, each being counted.

A word-form occurring in the lexicon is segmented according to it.
A word not occurring in the lexicon is first bounded by the longest suffix such that the resulting stem has at least three characters; then it is bounded from the left by the longest prefix such that the resulting stem has at least three characters, if the suffix and prefix have each been seen at least 20 times.
Finally we encode the subword units, using a free marker BPE tokenizer trained on the pre-segmented text so as to not enable a merge crossing a Lemma2Word boundary.
The variant called \emph{lexicon} does not use the lemmatizer to obtain the lexical units, but does use it to construct the lexicon.

The \emph{online} variant uses Stanza's lemmatizer to obtain the lemma of the word being tokenized, and then uses the lemma to obtain the stem, as above; if the word and the lemma have no common substring, we fall back to the lexicon rule.
This online variant is the same as in the original report.
Unlike the lexicon variant, the online variant requires a tagger at tokenization time, and it shares training data with the gold standard, as shown in Section~\ref{sec:gold}.

\subsection{Reference Tokenizers}
We implemented two reference tokenizers.
One is a BPE tokenizer with a word-initial marker and the Metaspace pre-tokenizer, as implemented in the Hugging Face \texttt{tokenizers} library.
The other is the unigram model tokenizer implemented in SentencePiece~\cite{kudo2018subword}, which uses all characters and does not perform Unicode normalization, which ensures that all tokenizers have the same character set.

\subsection{Using a Language Model as a Segmenter}\label{sec:llm}
In the original report, we explored the possibility of utilizing a large language model to segment words.
We assess it as a segmenter rather than a tokenizer, since we provide the model with a fixed set of words to evaluate and the prompt instructs the model to mark morpheme boundaries without modifying any characters.
These boundaries are subsequently assessed in the same way as other methods.
The findings of this study are discussed in Section~\ref{sec:res-llm}, and the prompt is detailed in Appendix~\ref{app:prompt}.
